\section{Translation with quantifiers}
\label{sec:fol-semantics-03}

Every first-order translation must state a signature and the intended range of
the variables.  Restricted quantification is an abbreviation:
\[
 \forall x\in A\,\varphi\quad\text{means}\quad
 \forall x\,(A(x)\to\varphi),
\]
and
\[
 \exists x\in A\,\varphi\quad\text{means}\quad
 \exists x\,(A(x)\land\varphi).
\]

\begin{example}[Quantifier order]
Let the variables range over people, with unary relation symbols
\(S\) and \(P\) for student and professor, and binary relation symbol \(A\)
for admires.  Then
\[
 \forall x\bigl(S(x)\to
   \exists y(P(y)\land A(x,y))\bigr)
\]
says that every student admires some professor.  The formula
\[
 \exists y\bigl(P(y)\land
   \forall x(S(x)\to A(x,y))\bigr)
\]
says that one professor is admired by every student.  The second formula
implies the first: its existentially quantified professor serves as a witness
for each student.  The first formula does not in general imply the second.
Quantifier order records dependence of a witness on earlier
universally quantified objects.
\end{example}

\begin{example}[Negated universal]
In the same signature, ``No professor admires every student'' is
\[
 \forall y\bigl(P(y)\to
   \neg\forall x(S(x)\to A(y,x))\bigr).
\]
By classical quantifier duality, this is equivalent to
\(\forall y(P(y)\to\exists x(S(x)\land\neg A(y,x)))\).  If there are no
students, the inner universal in the first formula is vacuously true and the
existential in the second formula is false; both formulas therefore require
that there be no professors.
\end{example}

